\documentclass[10pt,twocolumn]{article}

\usepackage[utf8]{inputenc}
\usepackage[T1]{fontenc}
\usepackage{lmodern}
\usepackage[margin=0.75in,top=1in,bottom=1in]{geometry}
\usepackage{amsmath,amssymb}
\usepackage{booktabs}
\usepackage{listings}
\usepackage{xcolor}
\usepackage{hyperref}
\usepackage{microtype}
\hypersetup{colorlinks=true,linkcolor=blue!60!black,citecolor=green!50!black,urlcolor=blue!60!black}

\lstset{basicstyle=\ttfamily\footnotesize,breaklines=true,frame=single,
        backgroundcolor=\color{gray!8}}

\usepackage{../shared/cerebrum-macros}
% Unified notation table for CEREBRUM papers
% Resolve naming conflicts: TSC uses \eta_T for temporal weight; CSA uses \eta for the same concept.
% All papers should import this file and use the macros below for consistency.

\newcommand{\etaT}{\eta_T}        % TSC temporal weight (renamed from generic \eta to avoid conflict)
\newcommand{\etaCSA}{\eta}        % CSA temporal decay weight (10-param formula, position 7)
\newcommand{\alphaCSA}{\alpha}    % CSA semantic similarity weight
\newcommand{\betaCSA}{\beta}      % CSA community score weight
\newcommand{\gammaCSA}{\gamma}    % CSA edge-type weight
\newcommand{\deltaCSA}{\delta}    % CSA distance penalty
\newcommand{\epsilonCSA}{\epsilon}% CSA hop decay
\newcommand{\zetaCSA}{\zeta}      % CSA PageRank prior weight
\newcommand{\iotaCSA}{\iota}      % CSA node recency weight
\newcommand{\muCSA}{\mu}          % CSA synthesis-density penalty
\newcommand{\thetaCSA}{\theta}    % CSA grounding confidence weight

% Graph notation
\newcommand{\KG}{\mathcal{G}}     % Knowledge graph
\newcommand{\Eset}{\mathcal{E}}   % Edge set
\newcommand{\Vset}{\mathcal{V}}   % Vertex/node set
\newcommand{\Cset}{\mathcal{C}}   % Community set
\newcommand{\csascore}{a(u,v,k)}  % CSA attention score

% Traversal notation
\newcommand{\beam}{B}             % Beam width
\newcommand{\hops}{H}             % Number of hops
\newcommand{\mrr}{\text{MRR}}     % Mean Reciprocal Rank
\newcommand{\hitatk}[1]{\text{H@#1}} % Hits@k


\title{\textbf{Production Knowledge Graph Reasoning:\\
Fault Tolerance, Streaming Ingestion, Metacognitive Maintenance,\\
and Structural Hole Patching}}

% Author block — shared across all CEREBRUM arXiv papers
\author{
  Bryan Alexander Buchorn \\
  Independent Researcher \\
  Las Vegas, NV, USA \\
  \texttt{bryan.buchorn@gmail.com}
}

\date{June 2026}

\begin{document}
\maketitle

% -----------------------------------------------------------------------
\begin{abstract}
We document the engineering lessons learned from deploying the \CEREBRUM{} Knowledge
Graph reasoning framework as a production service. Five failure classes emerge in
production that are absent in research prototypes: traversal failures mid-hop, persistence
write failures, stream interruptions, process spawning failures, and state loss across
restarts. We describe five corresponding fault-tolerance patterns that bound the blast
radius of each failure class. We further describe the REM Cycle — a background
metacognitive maintenance loop for graph pruning, hallucination prevention, and
community re-partitioning — and the \THALAMUS{} ingestion pipeline for real-time
streaming edge ingestion. Finally, we catalogue eight cross-feature interaction bugs
discovered during integration that are not visible from individual component tests.
The result is a systems paper about the gap between KG reasoning \emph{research} and
KG reasoning \emph{production}: the algorithms described in prior work are necessary
but not sufficient. The engineering described here is what makes them reliable.
\end{abstract}

% -----------------------------------------------------------------------
\section{Introduction}
\label{sec:intro}

The academic literature on Knowledge Graph reasoning focuses almost exclusively on
benchmark performance. Papers report H@1, MRR, and F1 on MetaQA, WebQSP, and Freebase.
What the literature does not report is what happens when:
\begin{itemize}
  \item A traversal crashes mid-hop, and the client needs \emph{some} answer, not none.
  \item A disk fills during a query logging write, and the exception propagates to the client.
  \item A community detection process spawning fails, and the reasoner has no partition.
  \item Ten months of accumulated Engram cache is lost on a pod restart.
  \item A stream client disconnects mid-reasoning and cannot tell if the server crashed.
\end{itemize}

These are not edge cases. They are the routine failure modes of any production service.
This paper documents five fault-tolerance patterns, two operational subsystems (REM
Cycle, \THALAMUS{} pipeline), and eight integration bugs that only appear when
components interact. We hope this systems paper serves as a reference for the gap between
KG reasoning research and KG reasoning production.

% -----------------------------------------------------------------------
\section{System Architecture}
\label{sec:arch}

The \CEREBRUM{} production stack consists of four operational layers:

\begin{enumerate}
  \item \textbf{\THALAMUS{}} — ingestion pipeline: entity normalization, STDP causal
        discretization, embedding, community detection.
  \item \textbf{\CORTEX{}} — reasoning engine: \CSA{} scoring, Bayesian Beam Search,
        answer extraction.
  \item \textbf{Persistence} — QueryLog (NDJSON append-only), Engram cache (JSON),
        graph state (GraphSnapshot).
  \item \textbf{API} — FastAPI REST server with streaming NDJSON endpoint, WebSocket
        telemetry bridge.
\end{enumerate}

The five fault-tolerance patterns target different points in this stack. We describe each
in terms of the failure class it addresses, the pattern applied, and the invariant it
preserves.

% -----------------------------------------------------------------------
\section{Five Fault-Tolerance Patterns}
\label{sec:patterns}

\subsection{Pattern 1: Partial-Result HTTP 200}

\textbf{Failure class:} Traversal exception mid-hop.

\textbf{Without pattern:} A traversal failure at hop 4 propagates to the route handler,
returning HTTP 500 with no answers — even though 3 hops of valid reasoning were completed.

\textbf{Pattern:} \texttt{BeamTraversal} checkpoints the best paths after each completed
hop into \texttt{\_partial\_paths}. If a later hop raises, the route handler catches the
exception and returns \texttt{\_partial\_paths} with HTTP 200. \texttt{QueryResponse}
gains two fields:
\begin{lstlisting}
partial: bool = False   # True if result is incomplete
error: Optional[str]    # Exception message for client logging
\end{lstlisting}

\textbf{Invariant:} A traversal failure never returns fewer results than the last
completed hop.

\subsection{Pattern 2: Write Failure Isolation}

\textbf{Failure class:} Disk-full or OOM during QueryLog/Engram writes.

\textbf{Without pattern:} \texttt{QueryLog.record()} raises \texttt{OSError}, propagating
to the route handler and converting a persistence failure into a query failure.

\textbf{Pattern:} All write calls in the query hot path are wrapped in
\texttt{try/except Exception}, logging the exception at WARNING and proceeding:
\begin{lstlisting}
try:
    query_log.record(seeds, answers, rel_seq)
except Exception as exc:
    logger.warning("QueryLog write failed: %s", exc)
\end{lstlisting}

The same pattern wraps the \texttt{GlobalRebalancer} background worker, preventing
a community detection exception from silently killing the rebalancer thread.

\textbf{Invariant:} Persistence failures never degrade query availability. The persistence
layer is a best-effort side channel.

\subsection{Pattern 3: Stream Error Signalling}

\textbf{Failure class:} Server-side exception during a streaming NDJSON query.

\textbf{Without pattern:} The stream client receives a truncated response with no
indication of whether the stream ended normally or abnormally.

\textbf{Pattern:} The streaming generator catches all exceptions and emits a terminal
error chunk before closing:
\begin{lstlisting}
{"type": "error", "message": "<exc>",
 "partial_results": [...], "hop_completed": 3}
\end{lstlisting}

Clients detect the error type by inspecting \texttt{response.type}. Partial results
are included so the client need not discard all work.

\textbf{Invariant:} A stream client can always distinguish normal completion from an
error, and always receives whatever partial results were computed before the error.

\subsection{Pattern 4: ProcessPool Sequential Fallback}

\textbf{Failure class:} Process spawning failure during community detection.

\textbf{Without pattern:} \texttt{best\_of\_n\_dscf()} uses a
\texttt{ProcessPoolExecutor} to run $n$ parallel DSCF/TSC attempts and return the
best-modularity partition. If process spawning fails (fork limit exceeded, Windows
restrictions, container runtime constraints), the entire community detection call raises.

\textbf{Pattern:} Process spawning errors are caught and the function falls back to
sequential execution in the calling thread:
\begin{lstlisting}
try:
    results = pool.map(run_dscf, seeds)
except (OSError, RuntimeError):
    logger.warning("Process pool failed; running sequentially")
    results = [run_dscf(s) for s in seeds]
\end{lstlisting}

\textbf{Invariant:} Community detection always produces a partition, even in constrained
execution environments. Sequential fallback is slower but correct.

\subsection{Pattern 5: Durable Engram with Two-Tier Startup}

\textbf{Failure class:} State loss of accumulated Engram relation-pattern cache across
pod restarts.

\textbf{Without pattern:} The Engram cache accumulates relation-pattern knowledge over
thousands of queries. A pod restart resets it to empty, discarding all accumulated
experience.

\textbf{Pattern:} \texttt{Engram.save(path)} serializes the cache to JSON after each
$k$-th query (configurable flush interval). On startup, the server follows a two-tier
recovery:
\begin{enumerate}
  \item Load \texttt{engram.json} if present (fast, in-memory restore).
  \item Replay \texttt{query\_log.ndjson} into the Engram if no snapshot exists
        (full replay, $O(q)$ where $q$ = query log length).
\end{enumerate}

\textbf{Invariant:} Accumulated reasoning experience survives pod restarts, bounded
only by the configured snapshot frequency.

% -----------------------------------------------------------------------
\section{The REM Cycle: Metacognitive Maintenance}
\label{sec:rem}

The fault-tolerance patterns prevent individual failures. The REM Cycle addresses a
different problem: long-running autonomous KG reasoners accumulate entropy. Speculative
STDP edges, unverified InsightEngine discoveries, and stale community partitions degrade
reasoning quality over time.

\subsection{Bilateral Verification}

An edge $E_{uv}$ with relation \texttt{INSIGHT\_LINK} is verified by:
\begin{equation}
\mathcal{T}(E_{uv}) = \mathbb{I}\bigl(\mathrm{inference\_score}(u,v) \geq \sigma\bigr)
\wedge \mathbb{I}\bigl(C(u) \approx C(v)\bigr)
\label{eq:triangulation}
\end{equation}

Edges failing $\mathcal{T}$ are candidates for pruning. The bilateral check ensures
that neither structural topology nor reasoning confidence alone is sufficient for
verification — both must agree.

\subsection{Recursive Hallucination Prevention}

Insight edges that are not independently validated receive a skeptical decay rate $\rho$:
\begin{equation}
c_{t+1} = c_t \cdot (\lambda \cdot \rho)^{\Delta t}, \quad \rho < 1.0
\label{eq:insight_decay}
\end{equation}

This prevents the reasoner from reinforcing its own unverified discoveries in a
positive feedback loop. Edges transition from \texttt{SPECULATIVE} to \texttt{GROUNDED}
only when they are validated by independent user queries.

\subsection{Community Re-partitioning}

The REM Cycle monitors modularity drift $\Delta Q_{total}$. When it exceeds a threshold,
the cycle spawns a background TSC re-run and performs an atomic community map swap:
\begin{lstlisting}
# Atomic swap — no query reads stale community map
with lock:
    engine.community_map = new_map
\end{lstlisting}

Post-swap, a \texttt{BridgeTwinEngine} hook prunes stale bridge records whose source
or destination community has changed, preventing zombie bridges from persisting after
re-partitioning.

\subsection{IKGWQ Evaluation}

The Incomplete KG with Question-and-Wondering (IKGWQ) benchmark tests reasoning quality
under controlled edge removal. The REM Cycle's Synaptic Bridge synthesis recovers
connectivity across removed edges:

\begin{table}[h]
\small\centering
\caption{IKGWQ benchmark (MetaQA): REM Synaptic Bridge synthesis.}
\label{tab:ikgwq}
\begin{tabular}{lcc}
\toprule
\textbf{Incompleteness Level} & \textbf{H@1 without REM} & \textbf{H@1 with REM} \\
\midrule
Level 0 (0\% removed)   & 46.1\% & 46.1\% \\
Level 2 (25\% removed)  & 31.2\% & 38.7\% \\
Level 4 (50\% removed)  & 19.5\% & 27.3\% \\
\midrule
Graceful Degradation AUC & — & \textbf{0.89} \\
\bottomrule
\end{tabular}
\end{table}

% -----------------------------------------------------------------------
\section{\THALAMUS{}: Production Streaming Ingestion}
\label{sec:thalamus}

The \THALAMUS{} ingestion pipeline processes streaming entity-relation-entity triples
with four production requirements:

\begin{enumerate}
  \item \textbf{Entity normalization}: deduplicate aliases (``NYC'' = ``New York City'').
  \item \textbf{Relation normalization}: standardize relation labels across sources.
  \item \textbf{Namespace isolation}: prefix entity IDs by source namespace
        (\texttt{text::}, \texttt{signal::}) to prevent cross-modal collisions.
  \item \textbf{Confidence and provenance}: tag each edge with ingestion source,
        timestamp, and confidence score at ingest time.
\end{enumerate}

The pipeline is composable: each stage can be enabled/disabled independently,
and custom stages can be inserted at any point. Default throughput on a single CPU core is $\approx$50,000 triples/second for pure
normalization (measured on AMD Ryzen 9 9950X3D, MetaQA KG: 43,234 entities / 124,680 edges);
STDP discretization adds $O(1)$ amortized overhead per triple via Lazy Decay.

% -----------------------------------------------------------------------
\section{Eight Cross-Feature Integration Bugs}
\label{sec:bugs}

The following bugs only manifest when multiple components interact. They are not visible
from unit tests of individual components.

\begin{enumerate}
  \item \textbf{Community swap mid-traversal.} A background TSC re-run swapping the
        community map mid-query causes the beam to use inconsistent community assignments
        across hops. \emph{Fix}: query snapshot isolation — community map is snapshotted
        at \texttt{traverse()} start and released at end.

  \item \textbf{Zombie bridges after re-partition.} Bridge Twin relay nodes created before
        a community re-partition reference stale community IDs. Post-rebalance, they
        point to non-existent or wrong communities. \emph{Fix}: post-rebalance hook
        in \texttt{BridgeTwinEngine} prunes all bridges with stale community assignments.

  \item \textbf{Engram replay loop.} On startup, replaying \texttt{query\_log.ndjson}
        into Engram triggers new query log writes (recording the replay queries), which
        grow the log indefinitely. \emph{Fix}: replay is tagged with
        \texttt{source="replay"} and suppressed from being logged.

  \item \textbf{Write-failure isolation masking OOM.} Pattern 2's broad
        \texttt{except Exception} catches \texttt{MemoryError} from OOM conditions,
        suppressing them as ``persistence failures'' when the real cause is insufficient
        heap. \emph{Fix}: re-raise \texttt{MemoryError} explicitly after logging.

  \item \textbf{Synthesis-density penalty double-counting.} When a path crosses two
        synthetic Synaptic Bridge edges, the $-\mu\sigma_d$ penalty in \CSA{}
        was applied twice, over-penalizing valid paths. \emph{Fix}: $\sigma_d$ is
        capped at 1.0 per path regardless of synthetic edge count.

  \item \textbf{Counterfactual cache stale after graph update.} The counterfactual
        re-ranking cache is keyed by query ID, not graph state. After a REM Cycle prunes
        edges, cached counterfactuals reference pruned paths. \emph{Fix}: cache TTL
        reduced to 30s; explicit invalidation on REM prune events.

  \item \textbf{ProcessPool worker inheriting open file handles.} On POSIX systems,
        forked DSCF workers inherit the parent's open \texttt{query\_log.ndjson}
        file handle. Concurrent writes from worker and parent corrupt the log.
        \emph{Fix}: log file handle is closed before spawning the pool and re-opened
        after pool completion.

  \item \textbf{Adaptive beam width racing community density probe.} The density probe
        for adaptive beam width reads community edge counts from the live adapter.
        If a bridge materialization completes between the probe and the traversal,
        the beam width is set for a lower-density graph than actually encountered.
        \emph{Fix}: density probe is snapshotted alongside the community map at
        traversal start.
\end{enumerate}

% -----------------------------------------------------------------------
\section{Operational Lessons}
\label{sec:lessons}

\paragraph{Design the failure path first.}
Each of the five fault-tolerance patterns was added after its failure class was
encountered in production. Designing the failure path before the happy path would have
prevented two of the eight integration bugs.

\paragraph{Unit tests are necessary but not sufficient.}
All eight integration bugs passed their respective component unit tests. Integration
tests that exercise component interactions are required for production confidence.

\paragraph{Persistence is a side channel, not a dependency.}
Pattern 2 embodies a general principle: persistence writes in the query hot path should
never prevent a query response. Treat them as best-effort side channels with explicit
failure boundaries.

\paragraph{Background workers need supervisors.}
The GlobalRebalancer thread and the TaskQueue worker both require exception supervisors
that log failures and keep the thread alive. Silent thread death is worse than logged
failure — at least logged failure is recoverable.

% -----------------------------------------------------------------------
\section{Conclusion}
\label{sec:conclusion}

Production KG reasoning requires engineering that is not described in the academic
literature. The five fault-tolerance patterns described here bound the blast radius of
the five most common failure classes encountered in production deployment. The REM Cycle
maintains long-term reasoning quality via metacognitive pruning and re-partitioning.
The \THALAMUS{} pipeline ingests streaming triples with entity normalization, namespace
isolation, and STDP causal discretization at production throughput. The eight
integration bugs catalogue the class of problems that only appear when components
interact — a class that unit tests cannot catch. We release all patterns, bugs, and
fixes as part of the open-source \CEREBRUM{} framework.

\section*{Acknowledgments}
The author gratefully acknowledges the use of Claude (Anthropic) as a research assistant
for mathematical formalization, code generation, and manuscript preparation. All
conceptual contributions, architectural decisions, experimental design, and intellectual
claims are solely the author's.

\bibliographystyle{plain}
\begin{thebibliography}{99}

\bibitem{lamport1978}
Lamport, L. (1978).
Time, clocks, and the ordering of events in a distributed system.
\emph{Communications of the ACM}, 21(7).

\bibitem{bernstein1983}
Bernstein, P.A., \& Goodman, N. (1983).
Multiversion concurrency control — theory and algorithms.
\emph{ACM Transactions on Database Systems}, 8(4).

\bibitem{diekelmann2010}
Diekelmann, S., \& Born, J. (2010).
The memory function of sleep.
\emph{Nature Reviews Neuroscience}, 11(2).

\bibitem{dean2004mapreduce}
Dean, J., \& Ghemawat, S. (2004).
MapReduce: Simplified data processing on large clusters.
\emph{OSDI}.

\bibitem{kleppmann2017}
Kleppmann, M. (2017).
\emph{Designing Data-Intensive Applications.}
O'Reilly Media.

\bibitem{buchorn2026flagship}
Buchorn, B.A. (2026).
CEREBRUM: Training-Free Knowledge Graph Reasoning via Community-Structured Graph Attention.
arXiv preprint.

\bibitem{buchorn2026report}
Buchorn, B.A. (2026).
CEREBRUM v2.71.0: Complete Technical Specification.
arXiv preprint [CEREBRUM\_REPORT\_ID].

\end{thebibliography}

\end{document}
